% 分水岭的一维教学例：梯度幅值剖面 f=(3,1,2,5,4,0,2,6,3,1,3)，位置 0 至 10。
% 上：三个局部极小各自形成盆地，坝在位置 3 与 7；下：只在位置 1、9 设标记时，
% 中间盆地经较低的鞍点（位置 3，高 5）并入左侧区域，坝只剩位置 7。宽约 96mm。
\begin{tikzpicture}[x=8mm,y=2.1mm,font=\figurefont\small]
\begin{scope}[yshift=0mm]
\fill[FigureBlueFill,draw=FigureBlue,line width=.6pt] (-0.42,0) rectangle (0.42,3);
\node[font=\figurefont\footnotesize] at (0,3.9) {A};
\fill[FigureBlueFill,draw=FigureBlue,line width=.6pt] (0.58,0) rectangle (1.42,1);
\node[font=\figurefont\footnotesize] at (1,1.9) {A};
\fill[FigureBlueFill,draw=FigureBlue,line width=.6pt] (1.58,0) rectangle (2.42,2);
\node[font=\figurefont\footnotesize] at (2,2.9) {A};
\fill[FigureStroke] (2.58,0) rectangle (3.42,5);
\fill[FigureYellowFill,draw=FigureYellow,line width=.6pt] (3.58,0) rectangle (4.42,4);
\node[font=\figurefont\footnotesize] at (4,4.9) {B};
\fill[FigureYellowFill,draw=FigureYellow,line width=.6pt] (4.58,0) rectangle (5.42,0);
\node[font=\figurefont\footnotesize] at (5,0.9) {B};
\fill[FigureYellowFill,draw=FigureYellow,line width=.6pt] (5.58,0) rectangle (6.42,2);
\node[font=\figurefont\footnotesize] at (6,2.9) {B};
\fill[FigureStroke] (6.58,0) rectangle (7.42,6);
\fill[FigureRedFill,draw=FigureRed,line width=.6pt] (7.58,0) rectangle (8.42,3);
\node[font=\figurefont\footnotesize] at (8,3.9) {C};
\fill[FigureRedFill,draw=FigureRed,line width=.6pt] (8.58,0) rectangle (9.42,1);
\node[font=\figurefont\footnotesize] at (9,1.9) {C};
\fill[FigureRedFill,draw=FigureRed,line width=.6pt] (9.58,0) rectangle (10.42,3);
\node[font=\figurefont\footnotesize] at (10,3.9) {C};
\draw[FigureStroke,line width=.7pt] (1,-1.2)--(1,-0.2);
\fill[FigureStroke] (1,-1.6) circle (0.7mm);
\draw[FigureStroke,line width=.7pt] (5,-1.2)--(5,-0.2);
\fill[FigureStroke] (5,-1.6) circle (0.7mm);
\draw[FigureStroke,line width=.7pt] (9,-1.2)--(9,-0.2);
\fill[FigureStroke] (9,-1.6) circle (0.7mm);
\draw[FigureStroke,line width=.6pt] (-.6,0)--(10.6,0);
\node[anchor=west] at (-.6,8) {(a) 无标记的浸没：三个盆地};
\end{scope}
\begin{scope}[yshift=-27.3mm]
\fill[FigureBlueFill,draw=FigureBlue,line width=.6pt] (-0.42,0) rectangle (0.42,3);
\node[font=\figurefont\footnotesize] at (0,3.9) {A};
\fill[FigureBlueFill,draw=FigureBlue,line width=.6pt] (0.58,0) rectangle (1.42,1);
\node[font=\figurefont\footnotesize] at (1,1.9) {A};
\fill[FigureBlueFill,draw=FigureBlue,line width=.6pt] (1.58,0) rectangle (2.42,2);
\node[font=\figurefont\footnotesize] at (2,2.9) {A};
\fill[FigureBlueFill,draw=FigureBlue,line width=.6pt] (2.58,0) rectangle (3.42,5);
\node[font=\figurefont\footnotesize] at (3,5.9) {A};
\fill[FigureBlueFill,draw=FigureBlue,line width=.6pt] (3.58,0) rectangle (4.42,4);
\node[font=\figurefont\footnotesize] at (4,4.9) {A};
\fill[FigureBlueFill,draw=FigureBlue,line width=.6pt] (4.58,0) rectangle (5.42,0);
\node[font=\figurefont\footnotesize] at (5,0.9) {A};
\fill[FigureBlueFill,draw=FigureBlue,line width=.6pt] (5.58,0) rectangle (6.42,2);
\node[font=\figurefont\footnotesize] at (6,2.9) {A};
\fill[FigureStroke] (6.58,0) rectangle (7.42,6);
\fill[FigureRedFill,draw=FigureRed,line width=.6pt] (7.58,0) rectangle (8.42,3);
\node[font=\figurefont\footnotesize] at (8,3.9) {C};
\fill[FigureRedFill,draw=FigureRed,line width=.6pt] (8.58,0) rectangle (9.42,1);
\node[font=\figurefont\footnotesize] at (9,1.9) {C};
\fill[FigureRedFill,draw=FigureRed,line width=.6pt] (9.58,0) rectangle (10.42,3);
\node[font=\figurefont\footnotesize] at (10,3.9) {C};
\draw[FigureStroke,line width=.7pt] (1,-1.2)--(1,-0.2);
\fill[FigureStroke] (1,-1.6) circle (0.7mm);
\draw[FigureStroke,line width=.7pt] (9,-1.2)--(9,-0.2);
\fill[FigureStroke] (9,-1.6) circle (0.7mm);
\draw[FigureStroke,line width=.6pt] (-.6,0)--(10.6,0);
\node[anchor=west] at (-.6,8) {(b) 以位置 1、9 为标记：两个区域};
\end{scope}
\node[font=\figurefont\footnotesize] at (0,-16.2) {0};
\node[font=\figurefont\footnotesize] at (1,-16.2) {1};
\node[font=\figurefont\footnotesize] at (2,-16.2) {2};
\node[font=\figurefont\footnotesize] at (3,-16.2) {3};
\node[font=\figurefont\footnotesize] at (4,-16.2) {4};
\node[font=\figurefont\footnotesize] at (5,-16.2) {5};
\node[font=\figurefont\footnotesize] at (6,-16.2) {6};
\node[font=\figurefont\footnotesize] at (7,-16.2) {7};
\node[font=\figurefont\footnotesize] at (8,-16.2) {8};
\node[font=\figurefont\footnotesize] at (9,-16.2) {9};
\node[font=\figurefont\footnotesize] at (10,-16.2) {10};
\end{tikzpicture}
